\section{Chinchilla y LLaMA: primero hay que precisar el objetivo de optimización}

Cuando un equipo decide entrenar su propio LLM, Chinchilla ofrecía la scaling recipe pública más clara de su momento. La conclusión clásica del artículo es que, con un compute de entrenamiento fijo, el número de parámetros del modelo y los tokens de entrenamiento deben escalarse de forma más equilibrada; su modelo representativo tiene unos 70B parámetros y 1.4T tokens. Lo que la clase busca corregir no es esta conclusión, sino su aplicación errónea a otro objetivo: \textbf{si el modelo se va a desplegar una y otra vez, ¿conviene sacrificar el costo de un entrenamiento que ocurre una sola vez para reducir el costo de inferencia a largo plazo?}

\subsection{Diferencia entre compute-optimal e inference-aware}

De forma aproximada, el cómputo del preentrenamiento de un dense Transformer puede escribirse como
\[
C_{\mathrm{train}}\approx 6ND,
\]
donde $N$ es el número de parámetros y $D$ es el número de tokens de entrenamiento. Elegir $N$ y $D$ con $C_{\mathrm{train}}$ fijo es el problema de Chinchilla; pero el objetivo de un producto a lo largo de su ciclo de vida se parece más a
\[
J(N,D)=C_{\mathrm{train}}(N,D)
+Q\cdot C_{\mathrm{infer}}(N),
\]
donde $Q$ es el volumen acumulado de generación durante el despliegue. Si $Q$ es muy grande, un modelo más pequeño, aunque se entrene con más tokens, puede reducir de forma notable el costo total y la latencia.

\lecturefigure{05-objective-mismatch.png}{La «Chinchilla trap» surge de confundir el objetivo de un presupuesto de entrenamiento fijo con el objetivo completo del producto.}{Subtítulos locales 09:20--11:40; Hoffmann et al. 2022; Touvron et al. 2023.}

\begin{knowledgebox}{Lectura de la figura: primero se confirma la restricción de presupuesto y después se discute qué es «óptimo»}
El óptimo de la izquierda significa «el presupuesto de este entrenamiento no puede aumentar»; el óptimo de la derecha significa «el entrenamiento ocurre una sola vez y la inferencia se repite». Ambos pueden ser razonables, pero no deben mezclarse. Lo que el curso llama over-train de un modelo pequeño consiste en seguir entrenando más allá del número óptimo de tokens para un compute fijo, no en repetir los mismos datos sin límite.
\end{knowledgebox}

El artículo de LLaMA reporta que los modelos de 6.7B y 13B se entrenaron con alrededor de 1.0T tokens, y los de 32.5B y 65.2B con alrededor de 1.4T tokens. Este diseño subraya que un modelo pequeño, con datos suficientes, puede alcanzar capacidades muy superiores a lo que sugería la intuición inicial, y que además es más fácil de desplegar para inferencia. No demuestra que «menos parámetros es mejor», sino que traslada la elección del modelo de un benchmark aislado a un problema conjunto de calidad, latencia, memoria de GPU, capacidad de procesamiento y escala de uso.

\begin{warningbox}{«Más tokens» no equivale automáticamente a un modelo más fuerte}
El beneficio de aumentar $D$ depende de la calidad de los datos, la tasa de repetición, la programación de la tasa de aprendizaje y la capacidad del modelo. Si los tokens nuevos son sobre todo datos repetidos de baja calidad, el entrenamiento puede limitarse a aumentar el costo; si la capacidad del modelo es insuficiente, el beneficio marginal también disminuye. El inference-aware training sigue requiriendo un scaling experiment, no basta con decir «hay que entrenar más tiempo el modelo pequeño».
\end{warningbox}

\teachervoice{El docente enfatiza que mucha gente usa «compute-optimal» como un elogio universal, sin aclarar si se optimiza la loss de entrenamiento, la calidad final, la latencia de inferencia o el costo total de propiedad. En una discusión de ingeniería, ante la palabra «optimal», la primera reacción debe ser preguntar por el objective, la constraint y el workload.}

\subsection{Resumen de la sección}

Chinchilla establece la regla de proporción para un presupuesto de entrenamiento fijo; la lección que la clase extrae de LLaMA es incluir la eficiencia de inferencia en el objetivo. La llamada Chinchilla trap es, en esencia, un desajuste de la función objetivo, no un error de la scaling law en sí.
